\section{Methods and Benchmark Design}
\label{sec:methods}

% Question: What does this section describe?
% Answer: NOVA and the injector as currently configured; settings, tolerances
% and the choices carried over from the August release are in the appendix.
% Evidence: The structure of this section; NOVA_ASTRA_EXECUTION_20260905_R1/
% STATUS.md. Settings and tolerances: appendices/appendix_implementation.tex.
This section describes the two methods of this project, NOVA
(Section~\ref{sec:nova}) and the injector (Section~\ref{sec:injector}), as
they are currently configured. NOVA grew out of an earlier release, frozen on
26 August 2026; the choices carried over from it, and the settings of both
methods, are collected in Appendix~\ref{app:implementation}.

\subsection{\nova{}}
\label{sec:nova}

% Question: What does NOVA estimate, and what does it hold fixed?
% Answer: Spectrum, limb darkening, continuum and background in one objective
% over detector samples; the response, curvature factor, background pattern
% and anchor, uncertainty scales and geometry are estimated beforehand.
% Evidence: Sections 1.3 and 1.4; config/MATCHED_COVARIANCE_ENSEMBLE_R1.json
% (estimator.rebuild).
The aim of NOVA is to infer the transmission spectrum from the detector
pixels themselves, so that the split between starlight and background is
decided in the same fit as the transit. A single forward model maps the
transmission spectrum $D(\lambda)$, the limb darkening $\bm{u}(\lambda)$, the
brightness of the star and the additive background to every retained detector
sample, and all four are estimated by minimising one objective. Other
ingredients are estimated beforehand from the same data and held fixed: the
spatial response, a curvature correction, the temporal pattern and expected
level of the background, the uncertainty scales, and the transit geometry,
which comes from a white-light fit, the one place where NOVA does extract
(Section~\ref{sec:own-white}). I run NOVA on the real
WASP-17 visit (720 integrations) and on the injected inputs of
Section~\ref{sec:injector} (229 integrations). The pixel layout, the off-trace pixel set,
the background maps, the per-order scales and the form of the curvature
correction are carried over from the earlier release
(Appendix~\ref{app:implementation}); every other fitted product, including
the geometry, is re-estimated on each input.

\subsubsection{The detector forward model}

% Question: How does NOVA map the physical quantities onto detector pixels?
% Answer: Continuum times exposure-averaged transit times normalised spatial
% response, times a curvature factor, plus an additive background.
% Evidence: proposals/NOMENCLATURE_PRESENTATION.md; proposals/
% handoff_nova_s_20260828/01_EXACT_NOVA_S_METHOD.md section 3;
% NOVA_SP_NOVA_S_SCIENCE_FREEZE_20260826_R1/code/sparse_varpro/
% fine_grid_transit_adapter.py and response_aware_varpro.py.
Within one detector column, the pixels of a trace see nearly the same
wavelength, and so NOVA organises the samples into groups, where a group $g$
contains the pixels $P_g$ of one order in one column. Denoting by $Y_{tp}$
the calibrated count rate of pixel $p$ in integration $t$, NOVA predicts
\begin{equation}
 \widehat Y_{tp}
 = \gamma_t\, C_{tg}\,
   \mathcal{T}_{tg}(D, \bm{u};\, \Omega)\,
   q^\star_{pg}
 + B_{tp},
 \label{eq:nova-forward}
\end{equation}
where $g$ is the group containing pixel $p$, $C_{tg}$ is the brightness of
the star in that group at time $t$ (the continuum), $\mathcal{T}_{tg}$ is the
transit factor, which dims this light according to the depth $D$ and limb
darkening $\bm{u}$ at the wavelengths of the group and the orbital state
$\Omega$, $\gamma_t$ is a small correction for slow curvature in the
brightness of the star, $q^\star_{pg}$ is the spatial response, which
distributes the light of the group over its pixels, with
$\sum_{p\in P_g} q^\star_{pg}=1$, and $B_{tp}$ is the additive background.
Only the stellar term is dimmed, so any light which the model assigns to the
background does not dim during the transit. In practice the transit is evaluated on
a fine wavelength grid and projected onto the pixels through the
line-spread kernel and the spatial response, following the public trace and
wavelength calibration \citep{BainesEtAl2023Trace,BainesEtAl2023Wavelength}
and its kernel \citep{DarveauBernierEtAl2022,JWSTPipeline2025}; the supports
of the two orders share no pixel.

\subsubsection{The shared spectrum and transit model}

% Question: How are the spectrum, orbit and limb darkening represented?
% Answer: 160 top-hat coefficients (158 free) in an exp/weighted-mean form,
% an exposure-averaged jaxoplanet transit at the white-light geometry, and a
% penalised smooth deviation from a STAGGER limb-darkening reference.
% Evidence: NOVA_SP_NOVA_S_SCIENCE_FREEZE_20260826_R1/reporting/
% v50_cycle2_reporting_inputs/shared_adaptive_observed_k160.json;
% configuration/variant_E12.json; reports/ONEOVERF_CAUSAL_INVESTIGATION_20260909.md
% (Delta); code/project_sources/nonlinear_detector_sp_v0p1_20260720_src/
% nonlinear_sp/limb_darkening.py (lines 9-12, 935-965); reports/
% NOVA_LIMB_DARKENING_WORDING_CORRECTION_20260923.md; public Ahsoka epf (period).
A planet has one transmission spectrum, whichever order records its light,
and so both orders share $D(\lambda)$. It is represented by $K=160$ top-hat
basis functions, about 0.01~$\mu\mathrm{m}$ wide below 2~$\mu\mathrm{m}$ and
about 0.04~$\mu\mathrm{m}$ above, whose coefficients are
\begin{equation}
  D_j = D_{\mathrm{global}}\,
  \frac{\exp\!\left[(\bm{H}_v\bm{v})_j\right]}
       {\sum_k w_k \exp\!\left[(\bm{H}_v\bm{v})_k\right]} ,
  \label{eq:depth-basis}
\end{equation}
where $D_{\mathrm{global}}$ is the achromatic depth, $\bm{v}$ are the free
shape coordinates, $\bm{H}_v$ maps them onto the subspace orthogonal to a
constant, and $w_k$ are fixed basis-width weights. This separates the overall
depth from the shape, whose weighted mean is zero, and keeps every depth
positive. Two red coefficients have no response support and are held
constant, which leaves 158 free depth coordinates. Otherwise the spectrum is
free: there is no atmospheric template, no smoothing between wavelengths, no
prior on its shape and no term which lets the depths of the two orders
differ.

The orbital state is
$\Omega=(P,\,a/R_\star,\,e,\,\omega_{\mathrm{peri}},\,b,\,t_0)$. The period
is that used by \citet{LouieEtAl2025} and the orbit is circular, whilst the
mid-transit time $t_0$, the scaled semi-major axis $a/R_\star$ and the impact
parameter $b$ come from NOVA's own white-light fit. Each integration lasts
about 44~s, during which the planet moves, so $\mathcal{T}_{tg}$ is the
transit averaged over the exposure, computed with jaxoplanet
\citep{Jaxoplanet2025}. The limb darkening follows the quadratic law
\citep{MandelAgol2002}, written in the coordinates $q_1,q_2\in[0,1]$ of
\citet{Kipping2013},
\begin{equation}
 u_1=2\sqrt{q_1}\,q_2, \qquad
 u_2=\sqrt{q_1}\,(1-2q_2).
 \label{eq:kipping-coordinates}
\end{equation}
Limb darkening and depth are partly degenerate, and so the spectral fit
allows only a smooth deviation from a theoretical reference computed from
STAGGER intensities \citep{MagicEtAl2015}: one offset and one slope in
normalised log-wavelength for each coordinate and order, eight in total,
penalised both for departing from the theory and for curvature in
wavelength. The white-light fit, by contrast, leaves the limb darkening free.
The injector uses a different stellar model, MPS-ATLAS
(Section~\ref{sec:per-read}), so this reference is independent of the limb
darkening of the injected signal.

\subsubsection{The empirical spatial response}

% Question: How is the spatial response estimated, and what is its known
% limitation?
% Answer: A Huber fit of each pixel to a depth-free transit template,
% normalised within groups and smoothed at fixed trace rank; the fixed
% integer rank causes centroid resets and a red-end amplitude excess.
% Evidence: products/original_v45_reference_qstar_r1_package/source/
% build_v47_empirical_template.py and aggregate_v47_qstar.py; products/
% original_A_group_calibration_adaptation_r1/build_nova_s_real_wasp17_responses_r1.py;
% NOVA-S handoff 01_EXACT_NOVA_S_METHOD.md section 6 (red-end limitation).
Because the starlight dims during the transit whilst the background does
not, the part of a pixel's time series which follows the transit measures how
much starlight falls on that pixel. I therefore build, for each order $o$, a
template $\phi_o(t)$ of the shape of the transit from the median behaviour of
its bright, stable pixels, scaled so that it carries no information about the
depth, and fit every pixel, with Huber weights, as
\begin{equation}
 Y_{tp} = A_{pg}\,\phi_{o}(t) + \ell_p + \kappa_p\,\tau_t + \varepsilon_{tp},
 \label{eq:qstar-fit}
\end{equation}
where $A_{pg}$ is the transit-correlated amplitude of the pixel, $\ell_p$ is
its level, $\kappa_p\tau_t$ is a linear drift in time $\tau_t$, and
$\varepsilon_{tp}$ is the residual. The response is the normalised amplitude,
$q^\star_{pg}=A_{pg}/\sum_{p'\in P_g}A_{p'g}$. Each amplitude comes from a dip of only
about 1.5\%, and so the amplitudes are smoothed across neighbouring columns at
a fixed integer position within the trace, with a strength chosen by
cross-validation between two halves of the integrations, which must also
agree in how they distribute the flux and where they place its centre. The
fixed integer position makes the centre of the response occasionally jump
between columns. In audits of the earlier release, this came with a response
at the faintest red columns which placed roughly 1--2\% more amplitude inside
the fitted pixels than the light delivered to them, and so made the transit
too shallow there. The effect has not yet been re-measured for each input.

\subsubsection{The continuum, the background and the noise}

% Question: How are the continuum, background, curvature factor and sample
% uncertainties modelled?
% Answer: Linear continuum per group; eight fixed maps with a constant and a
% G1 amplitude, softly anchored to an off-trace fit; a shared quadratic
% curvature factor with fold checks; ERR times order and pixel scales.
% Evidence: build_e13b_basis_products.py (cluster-only); products/
% original_A_group_calibration_adaptation_r1/build_nova_s_real_wasp17_e13b_r1.py
% and build_nova_s_real_wasp17_source_curvature_r1.py; source/
% run_nova_original_A_group_detector.py; reports/
% OFFTRACE_TRANSIT_COMPONENT_DIAGNOSIS_20260921.md section 2; config/
% CURVATURE_POLICY_AUTHORIZATION_20260919.json; NOVA_SP_NOVA_S_SCIENCE_FREEZE_
% 20260826_R1/code/runtime/N_SCALE_CONTRACT.json; reports/
% offtrace_transit_diagnosis_20260921/sources/v47_oot_guard.py.
The continuum of each group is linear in time,
$C_{tg}=\beta_{0g}+\beta_{1g}\tau_t$, and the background is a sum of eight
fixed spatial maps $\Phi_{pk}$, each with a constant amplitude $\alpha_{k0}$
and one, $\alpha_{k1}$, which follows a temporal pattern $G_1(t)$,
\begin{equation}
 B_{tp}=\sum_{k=1}^{8}\Phi_{pk}
 \left[\alpha_{k0}+\alpha_{k1}\,G_1(t)\right].
 \label{eq:background-model}
\end{equation}
Because a scaled zodiacal model has already been subtracted from the images
(Section~\ref{sec:detector-processing}), $B_{tp}$ describes what remains,
including any error in that scaling. The first map is a standardised
calibration background and the other seven are smooth polynomial surfaces.
Fitting the maps to the off-trace pixels of every integration gives a time
series of amplitudes, and $G_1$ is the first of its principal components whose
correlation with the shape of the transit is at most 0.25 in absolute value,
which limits how much transit-shaped variation the background can absorb. On the real visit, the leading
component itself follows the transit closely, which means that part of the
off-trace light does dim with the star. The criterion keeps this component
out of the background model, the uncertainty ensemble puts it back
(Section~\ref{sec:ensemble}), and the far-light family of the injector
is designed to bracket its effect (Section~\ref{sec:dimmed-light}). A joint fit to all
off-trace pixels gives an expected value $\bm{\alpha}_{\mathrm{off}}$ of the
sixteen amplitudes, with precision matrix $\bm{N}_{\mathrm{off}}$, which acts
as a soft anchor: the background is refitted at every step, but cannot wander
far from what the off-trace pixels show.

A linear continuum cannot follow a gentle curvature in the brightness of the
star, which could leak into the depth. The factor $\gamma_t$ is therefore the
ratio of a quadratic fit to the out-of-transit flux of each order's trace to
its straight-line approximation, normalised to unit mean out of transit, and one
factor is shared by both orders. It is only accepted if two alternating halves
of the out-of-transit integrations predict nearly the same, small correction
during the transit and, for the real visit and the injected inputs, are
strongly correlated; for the ensemble realisations the correlation was made
a diagnostic after it was found to reject harmless corrections of
1--15~ppm. Finally, each
sample has uncertainty $\sigma_{tp}=\mathrm{ERR}_{tp}\,N_{o}\,s_{p}$, where
$\mathrm{ERR}_{tp}$ is the pipeline uncertainty, $N_o$ is a scale for order
$o$ (3.07 and 4.19) and $s_p\geq1$ is a scale for each pixel, estimated by
predicting one half of the out-of-transit integrations from the other. These
scales only set the weights of the fit; the uncertainty of the real spectrum
comes from the ensemble.

\subsubsection{The objective and its minimisation}

% Question: What is minimised, and how?
% Answer: A Huber data term plus a background anchor and two limb-darkening
% penalties, minimised by variable projection inside bounded TRF inside
% IRLS, with strict acceptance and at least two starts.
% Evidence: NOVA_SP_NOVA_S_SCIENCE_FREEZE_20260826_R1/code/sparse_varpro/
% production_irls.py and code/project_sources/nonlinear_detector_sp_v0p3_
% huber_discrepancy_factorial_20260722_r4_src/nonlinear_sp_huber_factorial/
% robust.py (HUBER_DELTA 1.345); code/runtime/STRICT_IRLS_POLICY.json;
% products/w17_measured_production_recovery_20260923_v7_r1_package/recipes/
% closure/source/close_causal_reduced_starts.py; reports/
% PHYSICAL_ENDPOINT_THIRD_START_20260918.md.
Writing the standardised residual as
\begin{equation}
 r_{tp}=\frac{Y_{tp}-\widehat Y_{tp}
 (D,\bm{u},\bm{\beta},\bm{\alpha};\,q^\star,\Omega)}{\sigma_{tp}},
 \label{eq:residual}
\end{equation}
where the fitted parameters appear before the semicolon and the fixed inputs
after it, the NOVA estimate is
\begin{equation}
\begin{split}
 &(\widehat D,\widehat{\bm{u}},\widehat{\bm{\beta}},
  \widehat{\bm{\alpha}})
 = \arg\min\Big\{ \\
 &\quad\sum_{(t,p)\in\mathcal{M}}\rho_{1.345}(r_{tp}) \\
 &\quad+\tfrac{1}{2}(\bm{\alpha}-\bm{\alpha}_{\mathrm{off}})^{\mathsf T}
   \bm{N}_{\mathrm{off}}(\bm{\alpha}-\bm{\alpha}_{\mathrm{off}}) \\
 &\quad+\tfrac{1}{2}\|\bm{r}_{\mathrm{LD,theory}}(\bm{u})\|_2^2
  +\tfrac{1}{2}\|\bm{r}_{\mathrm{LD,curv}}(\bm{u})\|_2^2
 \Big\},
\end{split}
 \label{eq:objective}
\end{equation}
where $\mathcal{M}$ is the mask of retained samples, $\rho_{1.345}$ is the
Huber loss, which is quadratic for $|r_{tp}|\leq1.345$ and linear beyond, so
that outliers do not dominate the fit, and $\bm{r}_{\mathrm{LD,theory}}$ and
$\bm{r}_{\mathrm{LD,curv}}$ are the scaled deviations of the limb darkening
from theory and its curvatures in wavelength. There is no Gaussian-process
term for correlated noise.

For given depths and limb darkening, the prediction is linear in the
thousands of continuum and background coefficients, and so NOVA splits the
parameters into a nonlinear block $\bm{\theta}=(D,\bm{u})$ and a linear
block $(\bm{\beta},\bm{\alpha})$ and minimises in three nested loops. In the innermost loop, with $\bm{\theta}$ and the Huber weights $w_{tp}$
held fixed, variable projection \citep{GolubPereyra1973} solves exactly for
the linear block, i.e.\ the least-squares problem with weights
$w_{tp}/\sigma_{tp}^2$ together with the background anchor. This is tractable
because each continuum coefficient is coupled only to neighbouring columns
through the line-spread kernel, and the sixteen background coefficients are
handled through a Schur complement; the Jacobian with respect to
$\bm{\theta}$ is obtained by differentiating through this exact solution. In the middle loop, bounded trust-region reflective least
squares \citep{BranchColemanLi1999} minimises this profiled objective over
$\bm{\theta}$, and in the outer loop the Huber weights
$w_{tp}=\min(1,\,1.345/|r_{tp}|)$ \citep{Huber1964} are updated. A fit is strict
only if the weights, the objective and the spectrum have all settled,
first-order optimality is reached and a final cycle with recomputed weights
reproduces the solution. Every analysis is run from at least two starts, one
at the white-light depth with a flat shape and one with a fixed perturbation
of the shape and limb darkening (with a third, at their midpoint, in the pilot
of the injector), and the strict start with the lowest objective is selected;
agreement between starts is reported but not required, and so far they agree
to within about 0.01~ppm.

\subsubsection{White-light geometry inference}
\label{sec:own-white}

% Question: Where does NOVA's transit geometry come from, and when is it
% accepted?
% Answer: Its own joint two-order Eureka!/dynesty white-light fit with wide
% geometry priors, accepted by the registered v7 grazing-mass rule.
% Evidence: source/nova_own_white_eureka.py, source/nova_own_white_geometry.py,
% source/v7_grazing_statistic_20260923.py; evidence/
% group_ownwhite_native_sources_r1/white_light.py; config/
% V7_GEOMETRY_GATE_REGISTRATION_20260923.json; reports/V7_CALIBRATION_RESULT_20260923.md;
% reports/SIZED_WHITE_CONTINUATION_RESULTS_20260921.json; public Ahsoka epf files.
NOVA infers $t_0$, $a/R_\star$ and the inclination $i$, and hence
$b=(a/R_\star)\cos i$, from its own detector data, without using any fitted
value from another pipeline or from the injector. For each order, a
white-light curve is built by summing a profile-weighted (optimal) estimate of
the light in each column, using an out-of-transit profile from the same
input, and the two curves are fitted jointly with the light-curve stage of
Eureka! \citep{BellEtAl2022}, sharing $t_0$, $a/R_\star$ and $i$, whilst each
order has its own radius ratio $k_o$, quadratic limb darkening, linear
baseline and white-noise multiplier. The geometry priors are deliberately
wide ($t_0$ uniform over the observed window, $a/R_\star$ on $[1.01,\,100]$
and $i$ on $[0^\circ,\,90^\circ]$), and the others, which use no fit to these
data, are copied from the public Ahsoka configuration of
\citet{LouieEtAl2025}. The posterior is sampled with
the nested sampler dynesty \citep{Speagle2020}, because on inputs injected at
the raw level the ensemble sampler emcee \citep{ForemanMackeyEtAl2013} did not
reach 50 autocorrelation lengths; a later check found about 12\% of the
samples of each failed chain on grazing or non-transiting geometries.

Such geometries are the failure that matters here. For every posterior
sample, $b$ is computed from that sample's $a/R_\star$ and $i$, and the
sample counts as grazing or non-transiting if
\begin{equation}
 b > 1-\max(k_1,\,k_2),
 \label{eq:grazing}
\end{equation}
i.e.\ if the planet's disc is not wholly inside the stellar disc at
mid-transit. A fit passes if the upper end of the 95\% interval of the
weighted fraction of such samples is below 1\%, passes with the fraction
reported if that upper end lies between 1 and 5\%, and otherwise fails, and a failed fit
is not continued, reseeded, trimmed or replaced by another geometry. This
rule is not blind to the real data. It is the fifth revision of the
acceptance criterion, written after the earlier versions had been abandoned
and after it was known that the real posterior contains no grazing samples of
positive weight. I did, however, fix it before testing it on saved
posteriors, where it had to reject the three emcee chains which had failed
their length assessment and accept both those which had passed and the nested
posterior of the real visit. An accepted fit passes its whole-posterior weighted medians of
$t_0$, $a/R_\star$, $i$ and $k_o$ to the spectral fit. For the real visit,
however, the accepted NOVA spectrum, the reference of the ensemble, was
fitted at the geometry of an earlier accepted emcee fit, which the later
nested fit reproduces closely.

\subsubsection{The uncertainty ensemble}
\label{sec:ensemble}

% Question: How is the uncertainty of NOVA's spectrum estimated and tested?
% Answer: 352 full-pipeline realisations of the real visit with block-bootstrap
% real noise; 256 estimate a rank-128 finite-ensemble Student likelihood, 32
% held out test its coverage.
% Evidence: config/MATCHED_COVARIANCE_ENSEMBLE_R1.json; products/
% generator_v2_20260923_v7_r2_package (REGISTRY.json, generator_v2_draw_20260922.py,
% pilot_noise.py); source/finite_ensemble_likelihood.py; source/
% covariance_validation.py; config/COVARIANCE_ANALYSIS_PROTOCOL_R1.json;
% evidence/reviewer_plan_reporting_rank_20260920_r1.json (rank 128).
Because the covariance of a single fit is conditional on uncertain
calibration products, such as the response and the geometry, I estimate the uncertainty of NOVA's spectrum by repeating the whole
analysis on simulated realisations of the real visit,
\begin{equation}
 Y^{(n)}_{tp}=\widehat Y^{\mathrm{ref}}_{tp}+\bar r_p+a_p\,\psi_t
 +\mathrm{ERR}_{tp}\,z_{d_n(t),\,p},
 \label{eq:ensemble-draw}
\end{equation}
where $n$ labels the realisation, $\widehat Y^{\mathrm{ref}}$ is the accepted
prediction for the real visit (with the fitted mean background off the
trace), and $\bar r_p$ is the mean real out-of-transit residual of each
pixel, so that any static mismatch between model and data is also present. The term $a_p\psi_t$ puts back the transit-shaped
off-trace component which the background model keeps out:
$\psi_t=\mathcal{T}^{\mathrm{w}}_t-1$ is the transit decrement of the
accepted white-light fit, exactly zero outside the transit, and $a_p$ is the
projection of the component onto $\psi_t$ in each off-trace pixel, extended
under the trace by the background model. The last term adds real noise, the
real ERR array times a standardised out-of-transit residual $z$ of a donor
integration $d_n(t)$; donors are drawn in blocks of 64 consecutive
integrations within one out-of-transit interval, because real noise is
correlated in time.

The ensemble contains 352 realisations, each with a fixed seed and role: 256
estimate the covariance, 32 are held out to test it, and two further sets of
32 test the block length, reusing the held-out seeds with blocks of 128 and
192 integrations. Each is
analysed from scratch, with its own calibration products, white-light fit
and two strict starts. Because a failed realisation may differ systematically
from a successful one, none is replaced or dropped: the covariance is only
formed once all 288 estimation and held-out realisations have passed their
checks, and otherwise the ensemble is reported as incomplete. The error of a realisation, its spectrum minus that of the real visit, is
expressed in $N_{\mathrm{m}}=128$ modes, the rank of the map onto the 147
bins. The $N=256$ estimation errors give the
sample mean $\bm{\mu}$ and the unbiased covariance $\bm{S}$, without shrinkage
or jitter. Because $\bm{S}$ is itself estimated, a Gaussian likelihood would be
too narrow, and so for a new error $\bm{x}$ I use the multivariate Student
distribution of \citet{SellentinHeavens2016}, centred on $\bm{\mu}$, with
$N-N_{\mathrm{m}}$ degrees of freedom and scale matrix
\begin{equation}
 \bm{\Sigma}=\left(1+\frac{1}{N}\right)\frac{N-1}{N-N_{\mathrm{m}}}\,\bm{S},
 \label{eq:finite-ensemble}
\end{equation}
so that $(\bm{x}-\bm{\mu})^{\mathsf T}\bm{\Sigma}^{-1}(\bm{x}-\bm{\mu})/
N_{\mathrm{m}}$ follows an $F(N_{\mathrm{m}},\,N-N_{\mathrm{m}})$
distribution, the factor $1+1/N$ allowing for the estimated mean. The
likelihood fails, and is not used for retrieval, if significantly fewer than
95\% of the held-out errors fall inside its 95\% region;
the shape of the distribution of the estimation errors is also compared with
simulated Gaussian ensembles, as a warning rather than a condition.

\subsection{The injector}
\label{sec:injector}

% Question: Why inject into raw ramps, and how is the injector used?
% Answer: Raw-level injection lets every method run its own detector
% processing on identical frames; a pilot tests the injector, production
% recovers many atmospheres in six visits.
% Evidence: source/emit_original_A_group_segment.py, source/
% original_v45_group_signal.py, source/group_accumulated_transport.py;
% evidence/matrix_trip_bindings_r2.json.
Real observations come with no ground truth, and so I create one by writing
a transit with a known spectrum into the raw detector ramps (the uncal files)
of a real observation, before any calibration, so that every recovery method
can run its own detector processing, including its own correction of the
$1/f$ noise, on the same injected frames. In a pilot stage, a few injected
cases are recovered with NOVA to test the injector and NOVA's sensitivity to
its assumptions; in production, a larger set of atmospheres will be injected
into each of six visits and recovered by every method.
\label{sec:baseline}

% Question: Which real data carry the injected signal, and what are the
% controls?
% Answer: The first 229 raw integrations, all before the real transit; a
% null with nothing injected and a rate-level injection.
% Evidence: evidence/matrix_trip_bindings_r2.json (baseline_provenance);
% reports/PHYSICAL_MATCHED_NATIVE_NULL_20260918.md; products/
% w17_measured_raw_emission_20260923_v7_r3_prepared/source/
% emit_conditional_matched_rate_20260920.py.
The data which carry the signal, the baseline, are the first 229
integrations of the raw WASP-17 observation, each of eight nondestructive
reads about 5.5~s apart; the real transit, with buffer integrations on either
side, only begins at integration 230. Two controls are built from the same
data: the null, the baseline processed with nothing injected, and a
rate-level injection, in which the count change of the final read, divided by
its exposure time, is added to the rate images of the null. Comparing the two
injections shows what the group-level processing does to the signal,
together with the effect of adding it as an integration average, which a ramp
fit does not exactly return when the transit changes during the reads.

\subsubsection{Which light is dimmed}
\label{sec:dimmed-light}

% Question: Which light does the injected transit dim?
% Answer: A source map in three zones: a modelled core, measured wings, and
% a far region injected at a fraction m (0, 1, 0.5, or the measured m_x).
% Evidence: products/light_based_family_20260921_r1/FAMILY.json and
% FAR_PIXEL_ACCOUNTING.json; products/full_light_outer_shape_controls_20260916_r1/
% SOURCE_ESTIMATE.json; products/full_photon_source_projection_diagnostic_r1/
% SOURCE_PROJECTION.npz; products/measured_per_column_20260922_r1/REGISTRATION.json;
% reports/EVENING_FAR_LIGHT_AND_WHITE_CHECKS_20260921.md.
In a real transit only the light of the target dims, and so the injector
needs a map of where that light falls, the source, which is divided into
three zones by the distance $d$ of each pixel from the nearer trace. In the
core ($d\leq16$ pixels) it is a projection of the stellar spectrum through
the public ATOCA model of the traces \citep{DarveauBernierEtAl2022},
calibrated on the same data. Further out it is the measured out-of-transit
light which remains after background subtraction, estimated without
smoothing and made nonnegative whilst preserving the total of each band of
columns, and the wings ($16<d\leq80$) take this estimate as it stands. In the far region ($d>80$) it is not known how much of the light
belongs to the target and how much to the sky, and so it is injected at a
fraction $m$ of the measured field. The members $m=0$ and $m=1$ of this family
bracket the possibilities, and $m=0.5$ is only used in the delivery check. Pixels which
are not eligible for the measured estimate, for example those contaminated
by field stars identified in the F277W exposure, keep the modelled light of
the core projection. The outer zones are distributed in wavelength and between the orders by a
separate model of the outer light, and the core spectrum varies linearly in
time between two calibrated spectra.

A fourth member uses a measured far-light fraction. If all of the far-light
dimming came from starlight, the ratio of the far-light dip to the transit
depth would equal the stellar share of that light, since the sky does not
dim. On the real visit, the far pixels dim during the transit by 0.834\% on
average. Dividing the dip of each column by the mean fractional dimming of the
star in the white-light model, taking a running mean over 128 columns and bounding the
ratio to $[0,1]$ gives a multiplier $m_x$ for column $x$, with a median of
0.606; this member injects 61.5\% of the measured far light.
Because $m_x$ is measured after the pipeline's $1/f$ correction, which may
imprint on the far rows, it is treated as an upper bound on the stellar
fraction, pending a repeat of the measurement before that correction, and
not as a measurement of it.

\subsubsection{The per-read transit signal}
\label{sec:per-read}

% Question: How is the transit computed, and which geometry is injected?
% Answer: ExoTETHyS TRIP with MPS-ATLAS set 1 intensities, averaged over each
% nested read window with a first time moment; the published b with a
% rescaled a/R* and inclination.
% Evidence: products/original_A_causal_validation_trip_reference_r1_package/
% source/run_original_v45_group_trip.py; source/B_mixed_native_trip.py;
% evidence/matrix_trip_bindings_r2.json (geometry); config/
% COMPARATOR_SHARED_GEOMETRY_PRIORS_20260923_R1.json (t0 prior).
The visible fraction of the starlight at time $\tau$, $T(\lambda,\tau)$, is computed for a
radius ratio $\sqrt{D(\lambda)}$ with the TRIP routine of ExoTETHyS
\citep{MorelloEtAl2020}, which integrates the occulted light over thin
annuli of the stellar disc, using MPS-ATLAS (set~1) intensities
\citep{KostogryzEtAl2022,KostogryzEtAl2023} at the stellar parameters of
WASP-17. The injected orbit keeps the period, the circular orbit and the
impact parameter ($b\simeq0.339$) of the public control file of
\citet{LouieEtAl2025}, but $a/R_\star$ and the inclination are rescaled so
that the transit, which really lasts about 4.3~h, lasts about 1.74~h and fits
centred in the baseline. No recovery method is given the rescaled values. The recoveries do share the
period, eccentricity and argument of periastron of the injected orbit, and
the prior on $t_0$ of every white-light fit is centred on, and started at,
the middle of the observed window, which is the injected mid-transit time by
construction.

Each read records the charge accumulated since the detector was reset, and
so read $n$ of integration $t$ sees the transit averaged over the window from
the reset to the end of that read. The transit is therefore averaged over
these eight nested windows by Gauss--Legendre quadrature and stored in
tables which hold, for each integration, read and wavelength, the window mean
$\langle T\rangle_{tn}(\lambda)$ and the first time moment about the
integration midpoint $\tau_t$,
$\mu_{tn}(\lambda)=\langle(\tau-\tau_t)(T-1)\rangle_{tn}$, which is needed
because the source itself changes during an integration. The tables are
checked against independent integration averages, and integration~0, whose
reads extend beyond the time range of the stellar calibration, must carry no
transit signal. An atmospheric transit, whose deeper wavelengths are occulted
slightly earlier and longer than the grey signal of the delivery check (integrations 51 to 177),
affects integrations 50 to 178 and leaves the other 100 unaffected.

\subsubsection{Transport to raw counts}
\label{sec:transport}

% Question: How is the count change written into the raw ramps?
% Answer: Per-read projected decrement with a source-slope term, converted
% back through the flat, then through an affine bridge and Newton inversion
% of the pipeline's linearity polynomial, written as floats.
% Evidence: source/original_v45_group_signal.py; source/
% group_accumulated_transport.py; config/DEVELOPMENT_GROUP_TRANSPORT_R1.json;
% NOVA_SP_MULTIVISIT_SCENE_ROBUSTNESS_20260826_R1/source/
% stage3b_detector_transport_r2.py; stage3b_float_sci_persistence_adapter_r2.py.
For a source which varies linearly in time, the count change accumulated by
read $n$ of integration $t$ at wavelength $\lambda$, and its projection onto
pixel $p$, are
\begin{equation}
\begin{split}
 \Delta N_{tn}(\lambda) = {}& s_n\big[S_t(\lambda)\,
 \big(\langle T\rangle_{tn}(\lambda)-1\big) \\
 &\quad +\dot S(\lambda)\,\mu_{tn}(\lambda)\big], \\
 \Delta Y^{\mathrm{raw}}_{tnp} = {}& f_p
 \sum_{o}\big[\bm{A}_o\,\Delta N_{tn}\big]_p ,
\end{split}
 \label{eq:v45-injector}
\end{equation}
where $s_n=5.494\,n$~s is the accumulated exposure time of read $n$, $S_t$ is
the stellar spectrum at the midpoint of the integration, $\dot S$ is its rate
of change, $\bm{A}_o$ is the detector projection of order $o$, which combines
the spectral-response matrix of the ATOCA kernel, a wing correction and the
calibrations of the source map, $f_p$ is the effective flat-field coefficient
of the pipeline, which converts the loss in flat-fielded counts to counts before
flat-fielding, and $\Delta Y^{\mathrm{raw}}_{tnp}$ is the resulting change in the
linearised counts of pixel $p$. The core and outer components are projected
separately, each with its own spectrum, and summed. The change is zero on the
reference pixels and must be nonpositive everywhere, since a transit can only
remove light; otherwise the injection stops.

The raw ramp records counts before the linearity correction, and so the
change is passed through its inverse. The raw counts are first carried to the
input of the pipeline's linearity step by an affine bridge, since the steps
before it only add fixed offsets and the injection never changes the
reference pixels. Newton's method
then solves the linearity polynomial for the input whose linearised value
equals the original plus $\Delta Y^{\mathrm{raw}}$, and the inverse bridge
maps it back to the raw sample. The modified ramps are written as
floating-point values, so the change is not rounded to integer counts.
Saturated samples keep their original values, and the undelivered part of
the decrement is recorded; the observed noise, background and field stars are
left untouched.

\subsubsection{Detector processing and delivery checks}
\label{sec:detector-processing}

% Question: How are the injected ramps processed, and how is delivery
% checked?
% Answer: NOVA's JWST-pipeline/exoTEDRF chain with group-level 1/f; bit-
% identity and sign checks, and a grey secant test of the 1/f gain across
% the far-light family.
% Evidence: source/run_nova_original_A_group_detector.py (lines 55-100);
% reports/PHYSICAL_MATCHED_NATIVE_NULL_20260918.md (exoTEDRF 2.4.3); evidence/
% native_null_completion_20260918_r1.json; reports/
% LIGHT_BASED_GREY_DELIVERY_RESULT_20260921.md; reports/
% MEASURED_GREY_DELIVERY_RESULT_20260922.md; source/
% run_fixed_family_delivery_report_20260920.py.
NOVA's input is prepared with the same chain as the real data, from the
JWST calibration pipeline \citep{JWSTPipeline2025} and exoTEDRF
\citep{Radica2024}. Firstly, the pipeline's group-level steps run up
to and including the reference-pixel correction. Secondly, exoTEDRF removes
the $1/f$ noise at group level, temporarily subtracting a scaled model of the
zodiacal background so that the stripes can be estimated. Thirdly, the
linearity, jump-detection, ramp-fitting and gain steps give count rates.
Finally, Stage~2 of exoTEDRF subtracts the scaled background model, applies
the flat field and corrects bad pixels. The $1/f$ step masks the field stars
identified in the F277W exposure, and I apply no further scene mask. On the
injected inputs, the background and $1/f$ steps take the first and last 50
integrations as their out-of-transit reference, exactly those which the
injected transit leaves unaffected. The null follows the same chain, whilst
each comparison pipeline runs its own processing from the raw ramps.

Before an injection is used, I check that the files have the expected number
of reads, that every count change is nonpositive, and that every pixel
outside the injected support, including the reference pixels, is identical
to the input bit for bit. The family is then tested with a grey signal, a constant
depth of 16{,}000~ppm, injected under the three far-light assumptions. The
concern is that the $1/f$ correction, estimated partly from the far rows,
might respond to the far light non-linearly. If the fraction of the signal
which survives the step, its gain $g_m$, is linear in $m$, then $g_{0.5}$ lies
on the line between $g_0$ and $g_1$, and the check requires the departure
from this line, as an equivalent depth, to be below 10~ppm in every order and
column where the signal is measurable; the measured-fraction member is
checked in the same way. Passing only shows that
the signal is delivered consistently across the family; whether the far
light actually comes from the star remains open.

\subsubsection{Carrying the far-light uncertainty into the results}
\label{sec:envelope}

% Question: How does the unknown far-light share enter the results?
% Answer: As an envelope from the two bracketing raw-level pilot recoveries,
% which sets the wording (validated within scope or usable), never release.
% Evidence: products/light_based_family_20260921_r1/FAMILY.json
% (endpoint_criterion, chromatic.cases); reports/
% PHYSICAL_INJECTOR_PROTOCOL_20260918.md section 5; reports/
% PHYSICAL_EMISSION_RUNTIME_AND_ARMS_20260918.md; reports/
% SOURCE_MAP_A_DECISION_20260920_R2.json.
Because the stellar share of the far light is unknown, its effect is carried
into the results as an envelope. In the pilot, a reference atmosphere from the
ATMO grid \citep{GoyalEtAl2019}, rescaled to WASP-17~b, is injected under
both bracketing assumptions, at the raw and the rate level, and recovered
with NOVA, each case with its own calibration products, geometry and three
strict starts. The envelope is the
root-mean-square difference between the two raw-level recoveries over the
147 reporting bins. A second atmosphere, with different chemistry, is
recovered only at the raw level and at $m=1$, and does not enter the
envelope. If the envelope is 10~ppm or less, I describe results obtained with
the family as validated within the scope of the source family, in the
limited sense that NOVA's spectra under the two bracketing assumptions differ
by at most 10~ppm root-mean-square; otherwise I describe them as usable and
quote the envelope with them. No result is withheld because of it. The
envelope measures only NOVA's sensitivity to this assumption, not a
confidence bound, the accuracy of the core map, or the response of the
comparison pipelines. The threshold of 10~ppm
is a design choice, a budget for an intended claim about differences of
30~ppm between pipelines.

In production, each atmosphere will be recovered once, with the
measured-fraction member, and described according to the pilot envelope.
There are 24 atmospheres per visit from the same grid, twelve with local
condensation and twelve with rainout chemistry, each fitted with two strict
starts, in six visits: WASP-17, HAT-P-26, WASP-39, HAT-P-18, WASP-96 and
WASP-80. Each visit only enters production once its own white-light geometry has
been accepted and its own maps pass the checks applied to WASP-17: the core
map must not exceed the light observed in the aperture, the wing map must be
stable between two halves of the data, and the family must pass the grey
delivery check. Any exception is declared and reported with the results.

\subsubsection{Comparison pipelines}
\label{sec:comparison}

% Question: How are the comparison pipelines run fairly?
% Answer: Current public versions, end to end on the same raw frames with
% their own samplers and criteria; two declared departures (joint TS white
% fit with a replaced LD grid; NOVA's wide geometry priors for all).
% Evidence: reports/COMPARATOR_AUTHOR_PROTOCOL_AUTHORIZATION_20260923.md;
% config/COMPARATOR_SHARED_GEOMETRY_PRIORS_20260923_R1.json; reports/
% TRANSITSPECTROSCOPY_JOINT_WHITE_AUTHORIZATION_20260922.md; reports/
% TRANSITSPECTROSCOPY_REMAINING_CONFIGURATION_20260922.md; reports/
% TS_AUTHOR_WHITE_COMPLETE_LD_SUPPORT_FAILURE_20260923.md.
NOVA is compared with current public versions of the three pipelines of
\citet{LouieEtAl2025}: Ahsoka, \texttt{transitspectroscopy} 0.3.12
\citep{Espinoza2022TransitSpectroscopy} and exoTEDRF 2.5.0
\citep{Radica2024}, the current release of supreme-SPOON, whose runs therefore
do not exactly replicate the paper-era reduction. For a fair comparison, each
is run as its authors run it, end to end on the same
injected raw frames, with its own sampler, run length and acceptance
criteria: Ahsoka fits the two orders separately with \texttt{emcee} in
Eureka!, exoTEDRF fits a joint two-order white-light curve with its own
nested sampler, and \texttt{transitspectroscopy} extracts a box aperture and
fits with \texttt{juliet} \citep{EspinozaEtAl2019}. Two departures from the
published configurations are needed. Firstly, the white-light curves of
\texttt{transitspectroscopy}, fitted separately by \citet{LouieEtAl2025}, are
fitted jointly here, following the package author's own implementation, with
the settings which the paper leaves open fixed in advance. One of these, the
limb-darkening grid, had to be changed after the first run failed at its
spectral stage, because the STAGGER grid has no models at this star's
gravity above 6500~K for ExoTETHyS to interpolate; it was replaced by
MPS-ATLAS set~2 before any spectral estimate or score existed. This set comes
from the same model family as the injected intensities, so, unlike NOVA's
reference (the nearest complete STAGGER plane), that of
\texttt{transitspectroscopy} is closely related to the injected signal. Secondly, since the published priors describe the real planet (those of
Ahsoka place $t_0$ about 2.5~h after the end of the injected window), every
white-light fit uses NOVA's wide priors on $t_0$, $a/R_\star$ and $i$.

No pipeline receives the rescaled injected geometry or the fit of another
pipeline. The injector and NOVA do, however, share the public trace and
wavelength calibration and the ATOCA kernel, so an error in these would be
written into the injected data and assumed again in NOVA's recovery, and the
benchmark could not detect it. My own convergence statistics for each pipeline are reported for information
only, and failed runs are reported alongside the successful ones.

\subsubsection{Reporting and scoring}
\label{sec:reporting}

% Question: How are spectra compared and scored, and how is the truth kept
% out of the analysis?
% Answer: A common 147-bin R=100 grid, no filling, RMSE/bias/shape RMS, and
% truth revealed only after all 147 estimates are recorded.
% Evidence: products/w17_frozen_postclosure_score_20260923_r1_prepared/source/
% score_w17_frozen147.py, original_score_reference.py and
% original_truth_averaging.py; config/CASE_TEMPLATE_UNBOUND.json; source/
% run_aperture_terminal_covariance_r1.py; AGENTS.md (8 September decision).
All spectra are compared on a common grid at $R=100$, onto which NOVA's
coefficients are mapped, only for reporting, as the overlap-weighted mean in
each bin. The reddest
bins, beyond the validated range of the public wavelength calibration, are
excluded, which leaves 147 bins ending at $2.740\,\mu\mathrm{m}$. No estimate
is interpolated or filled in, and a method without an estimate in one of these
bins is recorded as a failure and is not scored. The injected depth in
each bin, $D_i^{\mathrm{inj}}$, is the unweighted average of $D(\lambda)$
over the bin. I score each recovery with the root-mean-square error and the
shape RMS,
\begin{equation}
 \mathrm{RMSE}=
 \left(\frac{1}{N_{\mathrm{b}}}\sum_i e_i^2\right)^{1/2},
 \label{eq:absolute-rmse}
\end{equation}
\begin{equation}
 \mathrm{RMS}_{\mathrm{shape}}=
 \left[\frac{1}{N_{\mathrm{b}}}\sum_i(e_i-\bar e)^2\right]^{1/2},
 \label{eq:morphology-rmse}
\end{equation}
where $e_i=\widehat D_i-D_i^{\mathrm{inj}}$ is the difference between the
recovered and the injected depth in retained bin $i$, $N_{\mathrm{b}}$ is the
number of retained bins, and $\bar e$ is the mean bias, which I also report,
together with the maximum absolute error. Since
$\mathrm{RMSE}^2=\bar e^{\,2}+\mathrm{RMS}_{\mathrm{shape}}^2$, the scores
separate a constant offset from an error in shape, and no offset correction is
applied to any reported spectrum. The same scores are given on six contiguous ranges and on two summary
ranges, below 2.1~$\mu\mathrm{m}$ and from 2.12~$\mu\mathrm{m}$ to the red end, the second
isolating the red end, where earlier comparisons differed most.

The injected spectrum is only compared with a recovery once that recovery's
estimates in all 147 bins have been recorded, so that it cannot influence any
choice of calibration, model or starting point; these depend only on the
data, the instrument reference files, the published priors and theory
references above, and the products carried over from the earlier release. Two
exceptions coincide with the design of the injection, although neither uses
the injected spectrum: the out-of-transit reference integrations of the
background and $1/f$ steps, and the centre of the prior on $t_0$. The uncertainty of NOVA's spectrum of the real visit is that of the
ensemble, once validated. Because the ensemble simulates the full visit, the
injected recoveries are instead reported with the covariance of their fit,
conditional on the calibration products, and with their residual and
convergence diagnostics; the comparison pipelines are reported with their own
uncertainties, which are not validated in the same way.
